\chapter*{全 Step 早見表}
\addcontentsline{toc}{chapter}{全 Step 早見表}
\markboth{全 Step 早見表}{}

どこに何が書いてあるかの一覧です。当日はここから引きます。

\begin{center}
\footnotesize
\begin{longtable}{@{}P{10mm}P{62mm}P{26mm}P{44mm}@{}}
\toprule
\hd{Step} & \hd{内容} & \hd{タグ} & \hd{いつ} \\
\midrule
\endhead
\multicolumn{4}{@{}l@{}}{\textbf{第 1 部　事前準備（前週準備まで）}}\\
1 & 体制とスケジュールを決める & \tALL & 2 か月前〜 \\
\multicolumn{4}{@{}l@{}}{\quad\textit{準備 1　Mulka2 のイベントの準備}}\\
2 & 機材と消耗品をそろえる & \tALL\ \tSI\ \tEMIT & 3 週間前 \\
3 & Mulka2 のファイルの置き場所を把握する & \tALL & \\
4 & イベントを新規作成する & \tALL & \\
5 & コースデータを取り込む（OCAD） & \tALL & コース確定後 \\
6 & クラスデータとスタート・フィニッシュ方式 & \tALL\ \tSI\ \tEMIT & \\
7 & 競技情報を直接上書きする & \tALL & 必要時 \\
8 & ステーション／ユニットの設定と時計合わせ & \tSI\ \tSIAC\ \tEMIT & 前週準備 \\
\multicolumn{4}{@{}l@{}}{\quad\textit{準備 2　エントリーリストの整形}}\\
9 & エントリーリストを整形する & \tALL & 締切後 \\
10 & ゼッケン番号を設計する & \tALL\ \tRANK & \\
11 & スタートレーン割とスタート時刻 & \tRANK\ \tALL & \\
12 & スタート順をシャッフルする & \tRANK & \\
13 & スタートリストを用途別に 4 種類作る & \tALL\ \tRANK & \\
\multicolumn{4}{@{}l@{}}{\quad\textit{合流}}\\
14 & レンタルカードを割り当てる & \tALL\ \tSI\ \tEMIT & \\
15 & カード番号の実機照合をする & \tALL\ \tSI\ \tEMIT & \\
16 & 当日の操作を通しで練習する & \tALL & \\
\midrule
\multicolumn{4}{@{}l@{}}{\textbf{第 2 部　前日準備}}\\
17 & 試走・ポ確でコースデータを検証する & \tALL\ \tSI\ \tEMIT & 前日まで \\
18 & バックアップ計時を準備する & \tALL & \\
19 & データをバックアップして共有する & \tALL & \\
20 & 印刷物をそろえる & \tALL\ \tRANK & \\
21 & 積み込みと最終確認 & \tALL & 前日 \\
\midrule
\multicolumn{4}{@{}l@{}}{\textbf{第 3 部　当日運営}}\\
22 & 起動順序を守って立ち上げる & \tALL & 開場前 \\
23 & ネットワークを組む & \tALL & 開場前 \\
24 & 時計を合わせる & \tALL\ \tSI & 開場前 \\
25 & ステーションを起動し，ポ確カードを読む & \tALL\ \tSI\ \tSIAC\ \tEMIT & 開場前 \\
26 & 当日申込を入力する & \tALL\ \tPRAC & 受付中 \\
27 & カード変更・クラス変更・代走・欠席 & \tALL & 受付中 \\
28 & 当日版のリストを配る & \tALL & 受付後 \\
29 & コース設定を当日に変更する & \tALL & 必要時 \\
30 & カードを読み取る & \tALL\ \tSI\ \tEMIT\ \tSIAC & 競技中 \\
31 & ペナチェックをする & \tALL & 競技中 \\
32 & 読めないカード・遅刻・未登録カード & \tALL & 競技中 \\
33 & 各種の印刷をする & \tALL & 競技中 \\
34 & レンタルカードを回収する & \tALL\ \tSI\ \tEMIT & 競技中 \\
35 & 表彰対象を報告する & \tALL\ \tRANK & 順位確定後 \\
36 & 未出走（DNS）入力と未帰還者の確認 & \tALL & 閉鎖前後 \\
37 & 当日トラブル対応一覧 & \tALL\ \tSI\ \tEMIT\ \tSIAC & 随時 \\
\midrule
\multicolumn{4}{@{}l@{}}{\textbf{第 4 部　事後処理}}\\
38 & 成績を確定させる & \tALL & 当日中 \\
39 & ラップと成績を公開する & \tALL & 翌日まで \\
40 & 公式成績表を作る & \tRANK & 1 週間以内 \\
41 & データを保存・アーカイブする & \tALL & 1 週間以内 \\
42 & 機材を整理して返却する & \tALL\ \tSI\ \tEMIT\ \tSIAC & 返却期日 \\
43 & 振り返りを残す & \tALL & 1 週間以内 \\
\bottomrule
\end{longtable}
\end{center}

\vspace{2mm}
\begin{Note}[第 5 部以降の使い方]
第 5 部は\textbf{大会の種類ごとに「上の標準とどこが違うか」}だけを書いた章です。
自分の大会の種類（ランキング大会／練習会／リレー／スプリント／複数日）を先に読んでから，
第 1 部に戻ると，やらなくてよい Step が分かります。

第 6 部のチェックリストは印刷用，第 7 部は逆引き索引，第 8 部は用語集です。

\textbf{JOY2Mulka を使う場合}の操作は，独立した章にせず，
該当する Step の中に\ \colorbox{j2mfr}{\color{white}\sffamily\bfseries\scriptsize\strut JOY2Mulka を使う場合}\
のボックスとして置いてあります
（Step 2・9・10・11・12・13・28，第 5 部 5.2）。
\end{Note}
